\newpage

\begin{multicols}{2}

\begin{table}[H]

\caption{Ejemplo de tabla con encabezados color azul}
\label{tab:ejemplo_tabla_azul}

\begin{threeparttable}

{\small\setstretch{1.35}\rowcolors{1}{Grey}{White}

\begin{tabular}{M{0.4\linewidth}M{0.4\linewidth}}\hline
    
    \sffamily\bfseries\cellcolor{Primario}\textcolor{White}{Encabezado 1} & \sffamily\bfseries\cellcolor{Primario}\textcolor{White}{Encabezado 2} \\ \hline
    
    Dato 1 & Dato 2\tnote{1} \\
    Dato 3 & Dato 4 \\
    Dato 5 & Dato 6 \\
    Dato 7 & Dato 8 \\

    \hline
\end{tabular}

}

\begin{tablenotes}\tiny % ex: \tnote{1} ... \item[1]
    \item[1] Ejemplo de nota a pie de tabla
\end{tablenotes}

\end{threeparttable}
\end{table}



\begin{table}[H]

\caption{Ejemplo de tabla con encabezados color blanco}
\label{tab:ejemplo_tabla_blanco}

\begin{threeparttable}

{\small\setstretch{1.35}\rowcolors{1}{Grey}{White}

\begin{tabular}{M{0.4\linewidth}M{0.4\linewidth}}\hline
    
    \sffamily\bfseries\cellcolor{White}\textcolor{Primario}{Encabezado 1} & \sffamily\bfseries\cellcolor{White}\textcolor{Primario}{Encabezado 1} \\ \hline
    
    Dato 1 & Dato 2\tnote{1} \\
    Dato 3 & Dato 4 \\
    Dato 5 & Dato 6 \\
    Dato 7 & Dato 8 \\

    \hline
\end{tabular}

}

\begin{tablenotes}\tiny % ex: \tnote{1} ... \item[1]
    \item[1] Ejemplo de nota a pie de tabla
\end{tablenotes}

\end{threeparttable}
\end{table}

\end{multicols}